\section{The nef summand on the algebraic-reduction fibers}
\label{sec:cycles}

Let $T$ be a smooth compact K\"ahler manifold of dimension $n$, and suppose
that $0<a(T)<n$.  After resolving its algebraic reduction, we have a
fibration $h:T\to W$ with $W$ smooth projective.  The objective of this
section is to show that a nef rational line bundle $M$ occurring in the
adjoint decomposition theorem has numerically trivial restriction to a
general fiber of $h$.  Proposition~\ref{desc:nef} will then descend its
class to $W$, after modifications.

The inductive decomposition on an individual fiber supplies a semiample
positive part only up to a flat twist.  To obtain sections on the total
space, we must choose the twists from a countable collection of global
line bundles.  Relative divisor cycles provide the required countable
collection of maps.  We first isolate the two ingredients of this
construction.

\subsection{Sections and coherent base change}

\begin{lemma}\label{cycle:sections}
Let $h:X\to W$ be a holomorphic algebraic reduction of a smooth connected
compact complex manifold, with connected fibers and smooth projective
base.  Let $(L_i)_{i\in I}$ be a countable collection of rational
holomorphic line bundles on $X$.  Outside a countable union of nowhere
dense subsets of $W$, a smooth fiber $F$ of $h$ satisfies
\[
                         \kappa(F,L_i|_F)\leq 0
                         \qquad\text{for every }i\in I.
\]
Any prescribed countable collection of further conditions holding on
dense open subsets of $W$ can be imposed on the same fiber.
\end{lemma}

\begin{proof}
Choose a positive integer $q_i$ such that $q_iL_i$ is an actual line
bundle.  For every positive multiple $m$ of $q_i$, the proper direct-image
theorem gives a coherent analytic sheaf
\[
                         \mathcal E_{i,m}=h_*\cO_X(mL_i)
\]
on $W$.  Work over the smooth-fibration locus, where the line bundle
is flat over the base.  On a dense open set where the dimension of
fiberwise sections is locally constant, coherent base change gives
the isomorphism
\begin{equation}\label{cycle:base-change}
 \mathcal E_{i,m}\otimes\C(w)
       \simeq H^0(X_w,\cO_{X_w}(mL_i|_{X_w}))
\end{equation}
by \cite[Section~7, Satz~5]{Grauert}.  It suffices here
to delete proper analytic subsets within dense open sets of $W$.
Fix an ample line bundle $H$ on $W$.  By GAGA and Serre's theorem,
$\mathcal E_{i,m}\otimes H^{k_{i,m}}$ is generated by global sections for
some integer $k_{i,m}$.  The projection formula and
\eqref{cycle:base-change} therefore show that all sections of
$mL_i|_{X_w}$ lift, after choosing a nonzero frame of $H^{k_{i,m}}$ at
$w$, to global sections of
\[
                    \cO_X(mL_i)\otimes h^*H^{k_{i,m}}.
\]

There are only countably many pairs $(i,m)$.  Choose $w$ in the smooth
locus of $h$, outside all the base-change exceptions and the further
generic exceptions in the statement.  If
$\kappa(X_w,L_i|_{X_w})>0$, some divisible $m$ has two sections whose
ratio is nonconstant on the connected fiber $X_w$.  Lifting these
sections as above gives a meromorphic function on $X$ whose restriction
to $X_w$ is that ratio.  Every meromorphic function on $X$ factors
through its algebraic reduction.  Such a function, when defined on a
nonempty open subset of the smooth fiber $X_w$, is constant there:
a local holomorphic section of $h$ through a point where the quotient
is holomorphic shows that the base function extends holomorphically
at $w$.  This contradicts the chosen nonconstant ratio.

The countable simultaneous choice is legitimate by Baire's theorem.
If a generic assertion is initially established on a dense open set,
a proper analytic exceptional subset of that open is still nowhere
dense in $W$.  Thus such generic assertions can also be included in the
choice.
\end{proof}

\subsection{Divisor incidence and a countable collection of fibrations}

We use the cycle-space results of Barlet, Lieberman, and Fujiki
\cite{Barlet,LiebermanCycles,Fujiki1978,Fujiki}.  In the form needed here,
the spaces of
effective cycles of a compact K\"ahler manifold have countably many
irreducible components; each component is compact and belongs to
Fujiki's class $\mathcal C$; see \cite[p.~189]{Fujiki}.  The locus of
cycles contained in a fiber of a proper map is an analytic relative
cycle space \cite[Proposition~3.2]{Fujiki1978}.  We also use the
meromorphic fiber map obtained by analytic flattening: after modifying
the base, the flat strict transform has an analytic family of fiber
cycles and hence a holomorphic map to the cycle space
\cite[Chapter~V]{Barlet}.  All parameter spaces and images in class $\mathcal C$
will be replaced by smooth compact K\"ahler models when needed.

The following observation explains why the numerical class of a divisor
can detect a semiample fibration without fixing its flat twist.

\begin{lemma}\label{cycle:divisors}
Let $F$ be a smooth connected compact K\"ahler manifold, let $A$ be a
rational line bundle on $F$, and suppose that a modification
$\rho:\widehat F\to F$ from a smooth compact K\"ahler manifold satisfies
\begin{equation}\label{cycle:fiber-decomposition}
 \rho^*A\num P+R,
 \qquad P\text{ semiample},
 \qquad R=N(c_1(\rho^*A))\text{ a rational divisor}.
\end{equation}
Let $g:\widehat F\to Q$ be the fibration defined by $P$, so that $Q$ is
normal projective and $P\qsim g^*H_Q$ for an ample rational line bundle
$H_Q$.  There is a dense open subset $U\subset F$, independent of the
divisors below, with the following properties for every sufficiently
divisible positive integer $m$.
\begin{enumerate}
\item There exists an analytic family of effective
divisors on $F$ with class $m c_1(A)$ that distinguishes the general fibers
of the meromorphic fibration $g\circ\rho^{-1}$.
\item For every effective divisor $D$ with class
$m c_1(A)$, membership
in $\operatorname{Supp}D$ is constant along the intersection of $U$
with a general smooth fiber of that meromorphic fibration.
\end{enumerate}
\end{lemma}

\begin{proof}
The difference between the two sides of
\eqref{cycle:fiber-decomposition} is flat as a rational line bundle.
By Lemma~\ref{pre:picard}, there is a flat rational line
bundle $\eta$ on $F$ such that
\[
                         \rho^*(A+\eta)\qsim P+R.
\]
Choose $m$ clearing the denominators and such that $mH_Q$ is very ample.
The divisors
\[
                   g^*D_Q+mR,
                   \qquad D_Q\in|mH_Q|,
\]
descend to an analytic family of divisors on $F$ with class $m c_1(A)$.
Indeed $\rho_*\cO_{\widehat F}=\cO_F$, so the projection formula
identifies sections of $m(A+\eta)$ with sections of its pullback.
On the isomorphism locus of $\rho$, away from $\operatorname{Supp}R$,
they distinguish general $g$-fibers, since hyperplanes distinguish
points of $Q$.

For an arbitrary effective divisor $D$ of this class,
Lemma~\ref{pre:negative}(1) gives
\[
                         \rho^*D\geq mR.
\]
The residual divisor $D'=\rho^*D-mR$ is effective and has class
$m c_1(P)$.  Let $G$ be a smooth compact connected fiber of $g$.  If
$D'$ does not contain $G$, its restriction is an effective divisor of
class zero on $G$, hence is empty: in positive dimension this follows
by integrating against a K\"ahler form to the appropriate power.
Consequently $\operatorname{Supp}D'$ either contains $G$ or misses it.
For a point fiber the same alternative is immediate.  Taking the image
of the isomorphism locus of $\rho$, and deleting
$\operatorname{Supp}R$ and the nonsmooth locus of $g$, gives the stated
open set $U$.  This choice uses only $\rho$, $R$, and $g$, and works for
all the divisors simultaneously.
\end{proof}

\begin{lemma}\label{cycle:detection}
Let $h:X\to W$ be a fibration of smooth connected compact K\"ahler
manifolds with $\dim W<\dim X$, and let $A$ be a rational line bundle
on $X$.  There is a countable collection of diagrams
\begin{equation}\label{cycle:diagrams}
 \begin{tikzpicture}[baseline=(current bounding box.center),>=Stealth]
 \node (Xi) at (0,1.15) {$X_i$};
 \node (Zi) at (2.7,1.15) {$Z_i$};
 \node (X) at (0,0) {$X$};
 \node (W) at (2.7,0) {$W$};
 \draw[->] (Xi) -- node[above] {$f_i$} (Zi);
 \draw[->] (Xi) -- node[left] {$p_i$} (X);
 \draw[->] (Zi) -- node[right] {$k_i$} (W);
 \draw[->] (X) -- node[below] {$h$} (W);
 \end{tikzpicture}
\end{equation}
where $p_i$ is a modification, $f_i$ is a fibration, and $X_i,Z_i$ are
smooth compact K\"ahler manifolds, with the following property.

For a very general smooth fiber $F=X_w$, suppose that a smooth compact
K\"ahler modification $\rho:\widehat F\to F$ has a decomposition
\eqref{cycle:fiber-decomposition} for $A|_F$, and that the fibration
$g:\widehat F\to Q$ defined by $P$ has positive-dimensional base.
For some $i$, the restriction of $f_i$ over $w$ is bimeromorphically
equivalent to $g$.  In particular, on a common resolution, their
general fibers coincide.  The collection and the exceptional subsets
of $W$ can be fixed before $F$, $\rho$, $P$, and $R$ are chosen.
\end{lemma}

\begin{proof}
We associate a map to each component of the relative divisor space by
recording all its divisors through a point.  On a fiber with the stated
decomposition, Lemma~\ref{cycle:divisors} makes these incidence supports
constant along the semiample fibration.  Integral multiplicities will
then give constancy of the cycle maps themselves, while the
distinguishing linear system will give the converse.

Put $d=\dim X-\dim W$, and let $W^\circ$ be the smooth-fibration locus
of $h$.  Consider the relative cycle spaces whose points are pairs
$(w,D)$, where $D$ is an effective $(d-1)$-cycle in $X_w$.  They are
closed subspaces of the product of $W$ with the corresponding compact
cycle components of $X$.  Thus they have countably many compact
irreducible components in class $\mathcal C$.

For each positive integer $m$ clearing the denominator of $A$, retain
the components meeting $W^\circ$ on which the divisor class is
$m c_1(A|_{X_w})$.  This is a condition on an entire component over
$W^\circ$: in a smooth family the cohomology class of an analytic
family of cycles is locally constant, and $c_1(A)$ supplies a global
section of the same cohomology local system.  After resolving a
component, its open part over $W^\circ$ is connected, so equality at
one point implies equality throughout that part.  Components not
dominating $W$ have proper analytic images; exclude all these images
from the eventual choice of $w$.

Fix a retained dominating component and a smooth compact K\"ahler
resolution $S$ of it.  Denote its divisor at $s\in S$ by $D_s$ and form
the incidence
\[
                    \mathcal I_S
                       =\{(x,s)\in X\times S:x\in\operatorname{Supp}D_s\}.
\]
We take the incidence of the analytic cycle family, including its
limits from the open part over $W^\circ$.  Over that open part it is a
family of pure $(d-1)$-dimensional cycles.  Hence every irreducible
component there dominates $S$ and has dimension $\dim S+d-1$.
Let $I_1,\ldots,I_t$ be the incidence components that dominate $X$.
The other components have proper analytic images in $X$; outside
those images they contribute nothing to incidence.

For each $I_j\to X$, analytic flattening gives a meromorphic map
recording its generic fiber as an effective cycle in $S$.  All these
cycles have dimension
\[
                         e=\dim S+d-1-\dim X.
\]
Taking the tuple of these maps, together with $h$, gives a meromorphic
map $\Phi_S$ from $X$ into $W$ times a finite product of compact cycle
components of $S$.  On a dense open subset of $X$, the union of the
supports of the recorded cycles is precisely
\begin{equation}\label{cycle:incidence-support}
                         \{s\in S:x\in\operatorname{Supp}D_s\}.
\end{equation}
In obtaining this open set we discard the images of the nondominating
incidence components and the exceptional sets of the fiber maps.
The target components and the image of $\Phi_S$ belong to class
$\mathcal C$.  Resolving the graph, taking Stein factorization, and
replacing the base by a smooth compact K\"ahler model therefore gives
a diagram \eqref{cycle:diagrams}.  The coordinate $h$ ensures that its
base maps to $W$.  If there is no dominating incidence component, we
retain only the coordinate $h$; such a component will not be needed
to detect a positive-dimensional $Q$.

This construction gives countably many diagrams.  Fix their models,
their incidence-recording open sets, and their Stein factorizations.
We may simultaneously require that their fibers over the eventual
$w$ be smooth, that $p_i$ restrict to a modification over $w$, and that
each of the dense open sets used meet the corresponding fiber densely.
For any one diagram these are generic conditions, by properness,
generic smoothness, and the fiber-dimension theorem.  In particular,
let $U_i\subset Z_i$ be the locus over which $f_i$ is smooth.  We
require that $k_i^{-1}(w)$ meet $U_i$ densely.  To justify this
condition, a component of $Z_i\setminus U_i$ that does not dominate
$W$ has proper image, whereas a component dominating $W$ has general
fiber of dimension strictly less than $\dim Z_i-\dim W$.
Since the general $k_i$-fiber is smooth of that pure dimension, it
cannot have a component contained in $Z_i\setminus U_i$.  The same
conditions may be imposed after any preselected countable collection
of further modifications of these diagrams.

We now choose such a very general $w$ and a decomposition on $F=X_w$
as in the statement.  By Lemma~\ref{cycle:divisors}, a projective
linear system of divisors of class $m c_1(A|_F)$ distinguishes the
general $g$-fibers for some $m$.  Its image in the relative cycle
space is irreducible and lies in one irreducible component.  That
component dominates $W$, since all images of nondominating components
were excluded.  It is therefore one of the components used above.
Passing to its parameter resolution $S$ preserves the distinguishing
property: every divisor still has a preimage in $S$.

The map $\Phi_S|_F$ distinguishes general $g$-fibers.  Indeed two
different general points of $Q$ are separated by a divisor in the
chosen linear system, so their incidence supports
\eqref{cycle:incidence-support} differ.  Conversely, fix a general
smooth connected $g$-fiber $G$.  On the common open set of
Lemma~\ref{cycle:divisors} and the incidence recording, membership in
every divisor $D_s$ is constant along $G$.  Thus the support in
\eqref{cycle:incidence-support} is fixed as the point moves in that
open part of $G$.

This support assertion also gives constancy of the recorded
\emph{cycles}.  Their fixed union of supports is pure $e$-dimensional
and has finitely many irreducible components.  Since all recorded
cycles have dimension $e$, there are only countably many possible tuples of
effective cycles supported there: their multiplicities are
nonnegative integers.  A holomorphic map from a connected complex
manifold with countable image is constant.  Removing a proper
analytic subset from the smooth connected $G$ leaves a connected
open set, so $\Phi_S$ is constant on a general $g$-fiber.  Point
fibers require no separate argument.

Constancy on the general connected fibers gives a meromorphic
factorization through $Q$.  One can see this by taking the image of
the graph in $Q$ times the target: its projection to $Q$ has one-point
general fibers and is bimeromorphic.  The distinguishing property
makes the induced map from $Q$ generically one-to-one onto its image.
Hence $\Phi_S|_F$ defines the same meromorphic fibration as $g$.

Finally, the restriction of $f_i$ over $w$ still has connected fibers:
its fibers are exactly the corresponding fibers of $f_i$.  Its base
over $w$ is connected, being the image of the connected source fiber.
The finite map in Stein factorization cannot split a connected
general $g$-fiber: a map from that fiber into a finite set is constant.
The preselected resolutions restrict to modifications over $w$.
Thus the same comparison holds for $f_i$, proving the assertion.
\end{proof}

\subsection{Vanishing of the nef class on the fibers}

We can now choose the flat twists on the total space.  The order of
choice is essential: the following proof fixes a countable collection
of global representatives before it chooses the smooth fiber.

\begin{proposition}\label{cycle:nef-summand}
Assume Theorem~\ref{thm:decomposition} in dimensions less than $n$.
Let $T$ be a smooth connected compact K\"ahler $n$-fold with
$0<a(T)<n$, and let $h:T\to W$ be a holomorphic algebraic reduction
with smooth projective base and connected fibers.  Let $B$ be an
effective rational SNC divisor with coefficients less than one,
assume that $J=K_T+B$ is pseudo-effective, and let $M$ be a nef
rational line bundle.  Then
\[
                              c_1(M|_F)=0
\]
for a general smooth fiber $F$ of $h$.
\end{proposition}

\begin{proof}
Set $A=J+M$.  By Theorem~\ref{pre:ordinary}, fix an effective rational
divisor $E$ with $J\qsim E$.  Apply Lemma~\ref{cycle:detection} to $h$
and $A$, and
fix its countable collection $(p_i,f_i,k_i)$.  Let $I_0$ be the subset
of indices for which $c_1(p_i^*M)$ vanishes on the general smooth
fiber of $f_i$.  For every $i\in I_0$, apply
Corollary~\ref{desc:representative} to $p_i^*M$ on $X_i$.  It gives a
representative on $X_i$ differing from $p_i^*M$ by a flat rational
line.  By Lemma~\ref{pre:picard}(2), that flat difference comes from
$T$.  We therefore obtain a rational line bundle $M_i^*$ on $T$, with
\[
                              M_i^*\num M,
\]
whose pullback is an actual rational pullback from the intermediate
base, after further modifications of the diagram.  Fix one such
representative and one such diagram for every $i\in I_0$.  In
particular, $M_i^*$ restricts rationally trivially to a general fiber
of that modified fibration.  This is a countable collection of
choices on $T$.

Choose $w\in W$ simultaneously satisfying the detection lemma and
its generic conditions for all the further diagrams just fixed.
Require also the conclusion of Lemma~\ref{cycle:sections} for the
countable collection
\[
                              (J+M_i^*)_{i\in I_0}.
\]
We take $F=T_w$ smooth, with $B|_F$ an SNC klt boundary.
We also require that the defining section of $E$ restrict
nontrivially to $F$; then
\[
                       J_F:=K_F+B|_F\qsim E_F:=E|_F\geq0.
\]
All the choices preceding $w$ depend only on data on the total space.

Suppose, towards a contradiction, that $M|_F\not\num0$.  The
lower-dimensional decomposition theorem gives a modification
$\rho:\widehat F\to F$ and
\begin{equation}\label{cycle:inductive-decomposition}
 \rho^*(J_F+M|_F)\num P_F+R_F,
 \qquad R_F=N(c_1(\rho^*(J_F+M|_F))),
\end{equation}
where $P_F$ is semiample and $R_F$ is a rational effective divisor.
Lemma~\ref{pre:negative}, applied to the nef summand $\rho^*(M|_F)$,
gives
\begin{equation}\label{cycle:residual-effective}
 R_F\leq N(c_1(\rho^*J_F))\leq\rho^*E_F,
 \qquad
 P_F\num(\rho^*E_F-R_F)+\rho^*(M|_F).
\end{equation}
Let $g:\widehat F\to Q$ be the fibration defined by $P_F$.  If
$Q$ were a point, the last equality would make a nef class and an
effective divisor sum to zero.  Pairing with a K\"ahler power would
give $\rho^*(M|_F)\num0$, contrary to the assumption.  Hence
$\dim Q>0$.

On a general smooth $g$-fiber $G$, the class of $P_F$ is zero.
Restricting the last equality in
\eqref{cycle:residual-effective} and using the same positivity
argument shows that
\begin{equation}\label{cycle:restriction-zero}
 c_1(\rho^*(M|_F)|_G)=0,
 \qquad
 G\cap\operatorname{Supp}(\rho^*E_F-R_F)=\varnothing.
\end{equation}
If $G$ is a point, the first assertion is automatic and the second
holds for a general point.  In positive dimension, choose a general
fiber not contained in the effective divisor, restrict it, and
integrate the effective and nef classes against a K\"ahler power.
Both nonnegative masses must vanish.

Lemma~\ref{cycle:detection} gives an index $i$ such that the restriction
of $f_i$ over $w$ and $g$ have the same general fibers up to modifications.  By
\eqref{cycle:restriction-zero} and injectivity of pullback on
cohomology, $p_i^*M$ has class zero on a smooth $f_i$-fiber above
$w$.  Such a fiber can be chosen in the smooth-fibration locus of
$f_i$, by the generic conditions already imposed.  Over the
connected smooth locus in $Z_i$, the restrictions of this Chern
class form a locally constant section of $R^2(f_i)_*\R$.  Vanishing
at one fiber therefore implies vanishing at every fiber in this
locus.  Thus $i\in I_0$; its representative $M_i^*$ was already
chosen before $w$.

We claim that
\begin{equation}\label{cycle:positive-kappa}
                         \kappa(F,J_F+M_i^*|_F)>0.
\end{equation}
On a common resolution of the maps on $F$, the line $M_i^*|_F$
is rationally trivial on a general $g$-fiber because it is a
pullback from the intermediate base of $f_i$.  This triviality
descends between smooth models of the fiber: for a modification,
pullback on line bundles is injective by normality and the
projection formula.  The difference
\[
             \Lambda=\rho^*(J_F+M_i^*|_F)-P_F-R_F
\]
is a flat rational line bundle by
\eqref{cycle:inductive-decomposition}.  Its restriction to a
general $G$ is rationally trivial.  Indeed $P_F|_G$ is trivial,
$M_i^*|_G$ is trivial, and
\eqref{cycle:restriction-zero} makes
$\cO_{\widehat F}(\rho^*E_F-R_F)|_G$ trivial, while
$J_F\qsim E_F$.

Resolve $Q$ and the induced map, and continue to denote the
resulting fibration between smooth models by $g$.  By Lemma~\ref{desc:flat},
$\Lambda\qsim g^*\lambda$ for a flat rational line bundle
$\lambda$ on the smooth projective base.  The line $P_F$ is the
pullback of the ample rational line on the original $Q$; on its
resolution the corresponding base line $H_Q'$ is nef and big.
Twisting by $\lambda$ preserves bigness.  On these models we have
an actual rational identity
\[
                \rho^*(J_F+M_i^*|_F)
                        \qsim g^*(H_Q'+\lambda)+R_F.
\]
Multiplication by the canonical section of a divisible multiple
of $R_F$ embeds the section spaces of the big base line into the
section spaces on $\widehat F$.  Since $\dim Q>0$, this proves
\eqref{cycle:positive-kappa}.

This contradicts the instance of Lemma~\ref{cycle:sections} imposed
on $J+M_i^*$ when $w$ was chosen.  Hence $M|_F\num0$ for the
chosen fiber.  Finally, restrictions of $c_1(M)$ form a
locally constant section of $R^2h_*\R$ on the connected smooth
locus of $h$.  Their vanishing therefore holds on every fiber in
that locus, and in particular on a general smooth fiber.
\end{proof}

\begin{corollary}\label{cycle:base-descent}
Under the assumptions of Proposition~\ref{cycle:nef-summand},
there are modifications $p:T'\to T$ and $b:W'\to W$, with
$T'$ smooth compact K\"ahler and $W'$ smooth projective, a
fibration $h':T'\to W'$ lifting $h$, and a nef rational line
bundle $N$ on $W'$ such that
\[
                              p^*M\num(h')^*N.
\]
\end{corollary}

\begin{proof}
Proposition~\ref{cycle:nef-summand} supplies the fiberwise
vanishing required in Proposition~\ref{desc:nef}.  Apply that
proposition to $h$ and $M$.  The modified base remains projective,
so the descended rational class is nef.
\end{proof}
